% Experimental setup. Sources: configs/inference.yaml, configs/activations.yaml,
% src/extract/{qwen,llava,cache_chartqa}.py, src/probes/*.py.
\section{Experimental Setup}
\label{sec:setup}

\paragraph{Models and inference.} We use two frozen VLMs, \qwen
\citep{bai2025qwen25vl} and LLaVA-NeXT with a Mistral-7B language model
\citep{liu2024llavanext}. Each question gets the standard ChartQA suffix
``Answer the question using a single word or phrase.'' and is decoded greedily
(at most 32 tokens) in fp16 on RTX 2080 Ti GPUs. Answers are scored by ChartQA
relaxed accuracy: numbers within 5\% of the gold, strings by exact match.

\paragraph{Activations.} Forward hooks record the residual stream after every
decoder layer during the prefill of the same \texttt{generate()} call that
produces the answer, so each vector belongs to the scored answer. We pool four
positions per layer (mean and max over image tokens, the last prompt token, and
the mean over prompt text tokens) for all 28 layers of Qwen ($d=3584$) and 32 of
LLaVA ($d=4096$). The image precedes the question and attention is causal, so
image-token states cannot depend on the question.

\paragraph{Probes.} A probe is an $L_2$-regularised logistic regression on one
pooled vector, z-scored with training-set statistics, with $\lambda$ chosen from
$\{1,10,10^2,10^3,10^4\}$. The binary probe uses a 60/20/20 split by
\emph{figure}: layer, position and $\lambda$ are chosen on validation and the test
set is scored once. Type probes, whose classes are smaller, use five folds by
figure: within each outer fold, layer and $\lambda$ are chosen by inner
cross-validation on the other folds, so every item is scored by a probe that
never saw its chart; the position is fixed to the last prompt token. We report
AUROC with 95\% intervals from a bootstrap over figures (2,000 draws), since
questions on one chart are correlated. Seed 42 throughout.

\paragraph{Baselines and controls.} Every probe is reported beside a baseline
without activations. On ChartQA this is the same probe on four surface features
(question source, question length, number of image tokens, numeric gold). We
also \emph{residualise}: regress the surface features out of the activations,
fitted on training rows only, and retrain \citep{sahoo2026linearprobes}. On
synthetic charts the baselines are a question-text classifier, the image-token
probe, and the probe on chart metadata (pilot, axis range, tick density, number
of bars, lookalike name, phrasing).
